\documentclass[11pt]{article}
\usepackage[margin=1.1in]{geometry}
\usepackage{amsmath,amssymb,amsthm}
\emergencystretch=1.2em

\newtheorem{theorem}{Theorem}
\newtheorem{lemma}{Lemma}
\newtheorem{proposition}{Proposition}
\newtheorem{corollary}{Corollary}
\theoremstyle{remark}
\newtheorem{remark}{Remark}

\newcommand{\R}{\mathbb{R}}
\newcommand{\E}{\mathbb{E}}
\newcommand{\tr}{\operatorname{tr}}
\newcommand{\diag}{\operatorname{diag}}
\newcommand{\Xt}{\widetilde{X}}
\newcommand{\Mh}{\widehat{M}}
\newcommand{\Lh}{\widehat{L}}
\newcommand{\HSIC}{\operatorname{HSIC}}

\title{Correctness of the Supervision-Dial Spectral Probe}
\author{}
\date{}

\begin{document}
\maketitle

\begin{abstract}
\noindent We prove that the supervision-dial spectral probe --- the estimator that
interpolates, through a single parameter $\lambda\in[0,1]$, between Laplacian-eigenmaps
manifold embedding ($\lambda=1$) and HSIC-maximizing supervised projection
($\lambda=0$) --- returns an \emph{exact global maximizer} of its objective at every
$\lambda$, that its two endpoints coincide with the established methods it claims to
interpolate, that its out-of-sample map is leakage-free under grouped cross-validation,
and that the solution path is well behaved in $\lambda$. All statements are
finite-sample and exact; statistical approximation guarantees (graph-Laplacian and HSIC
consistency) are inherited from the cited literature and are not re-proved here.
\end{abstract}

\section{Setup and the algorithm}

Fix a training fold $\{(x_i,y_i)\}_{i=1}^{m}$, $x_i\in\R^{p}$, $y_i\in\R$, and let
$P:\R^{p}\to\R^{r}$ be the PCA map fitted on the training fold only ($r\le m-1$),
followed by a fixed positive rescaling; write $\tilde x_i = P(x_i)$ and stack rows into
$\Xt\in\R^{m\times r}$. Define:

\begin{enumerate}
\item \textbf{Graph term.} $W\in\R^{m\times m}$: symmetric $k$-NN heat-kernel
  affinities, $w_{ij}=\exp(-\|\tilde x_i-\tilde x_j\|^2/\sigma_x^2)>0$ on retained
  edges, $w_{ij}=0$ otherwise, $w_{ii}=0$; degrees $d_i=\sum_j w_{ij}$,
  $D=\diag(d_1,\dots,d_m)$, Laplacian $L=D-W$.
\item \textbf{Dependence term.} $[K_y]_{ij}=\exp\!\big(-(y_i-y_j)^2/2\sigma_y^2\big)$;
  $H=I_m-\tfrac1m\mathbf{1}\mathbf{1}^{\!\top}$; $M=HK_yH$.
\item \textbf{Normalization.} $\Mh=M/\|M\|_2$, $\Lh=L/\|L\|_2$ (spectral norms;
  both matrices are nonzero in any nondegenerate fold).
\item \textbf{Pencil.} For $\lambda\in[0,1]$,
  \[
  A_\lambda \;=\; (1-\lambda)\,\Xt^{\!\top}\Mh\,\Xt \;-\; \lambda\,\Xt^{\!\top}\Lh\,\Xt,
  \qquad
  B \;=\; \Xt^{\!\top} D\,\Xt + \varepsilon I_r,\quad \varepsilon>0 .
  \]
\end{enumerate}

The \emph{probe} solves, for each $\lambda$ on a grid, the program
\begin{equation}\label{eq:program}
\max_{V\in\R^{r\times d}}\;\; J_\lambda(V) \;=\; \tr\!\big(V^{\!\top} A_\lambda V\big)
\qquad\text{s.t.}\quad V^{\!\top} B\, V = I_d ,
\end{equation}
by computing the $d$ eigenvectors of the generalized symmetric eigenproblem
$A_\lambda v=\gamma B v$ with largest eigenvalues, $B$-orthonormalized, discarding
directions whose embedding is constant across the training fold (Proposition~\ref{prop:trivial});
the embedding is $Z=\Xt V$ on train and $z_\ast=V^{\!\top}P(x_\ast)$ on any new point.
Note $J_\lambda(V) = (1-\lambda)\tr(Z^{\!\top}\Mh Z)-\lambda\tr(Z^{\!\top}\Lh Z)$ for
$Z = \Xt V$: \eqref{eq:program} is the joint embedding objective restricted to linear
maps of the (train-PCA) features, which is what furnishes the out-of-sample extension
(Locality Preserving Projections; He \& Niyogi \cite{he2003lpp}).

\section{Well-posedness}

\begin{lemma}[The constraint metric is positive definite]\label{lem:pd}
$B\succ 0$. Consequently the pencil $(A_\lambda,B)$ has $r$ real eigenvalues
$\gamma_1(\lambda)\ge\dots\ge\gamma_r(\lambda)$ with a $B$-orthonormal eigenbasis, and
\eqref{eq:program} is a maximization of a continuous function over a compact set, so a
maximizer exists.
\end{lemma}

\begin{proof}
Each vertex of the $k$-NN graph ($k\ge1$) has at least one incident edge, and heat
weights on edges are strictly positive; hence $d_i>0$ for all $i$ and $D\succ0$, so
$\Xt^{\!\top}D\Xt\succeq0$ and $B\succeq\varepsilon I\succ0$. For $B\succ0$ the change of
variables $W=B^{1/2}V$ turns $(A_\lambda,B)$ into the symmetric matrix
$\tilde A_\lambda = B^{-1/2}A_\lambda B^{-1/2}$, whose spectral decomposition yields the
stated eigenpairs; the constraint set $\{V:V^{\!\top}BV=I_d\}$ is the preimage of the
compact Stiefel manifold under the linear homeomorphism $B^{1/2}$, hence compact.
\end{proof}

\begin{lemma}[Smoothness identity]\label{lem:dirichlet}
For any $Z\in\R^{m\times d}$ with rows $z_i$,
$\;\tr(Z^{\!\top}LZ)=\tfrac12\sum_{i,j} w_{ij}\|z_i-z_j\|^2\ge 0$; in particular
$L\succeq0$, and if the graph is connected then $\ker L=\operatorname{span}\{\mathbf 1\}$.
\end{lemma}

\begin{proof}
$\tr(Z^{\!\top}LZ)=\sum_a z^{(a)\top}Lz^{(a)}$ columnwise, and for a vector $u$,
$u^{\!\top}Lu=\sum_i d_i u_i^2-\sum_{i,j}w_{ij}u_iu_j
=\tfrac12\sum_{i,j}w_{ij}(u_i-u_j)^2$ by symmetry of $W$. Nonnegativity is immediate;
on a connected graph $u^{\!\top}Lu=0$ forces $u_i=u_j$ along every edge, hence $u$
constant.
\end{proof}

\begin{lemma}[Dependence identity and what it measures]\label{lem:hsic}
With the linear kernel on $Z$ and kernel $K_y$ on $y$, the (biased) empirical HSIC of
Gretton et al.\ \cite{gretton2005hsic} is
$\widehat{\HSIC}(Z,y)=(m-1)^{-2}\tr\!\big(Z^{\!\top}MZ\big)$, i.e.\ the first term of
$J_\lambda$ up to the positive constants absorbed in $\Mh$. Its population counterpart is
$\HSIC(Z,y)=\sum_{a=1}^{d}\big\|C_{z_a\varphi(y)}\big\|_{\mathrm{HS}}^2$, the summed
squared cross-covariance operators between each embedding coordinate and the Gaussian
RKHS $\mathcal H_y$ on $y$; it vanishes if and only if
$\E[Z\mid y]=\E[Z]$ almost surely (mean-independence of the embedding from the year).
The empirical estimator converges at rate $O(m^{-1/2})$
\cite[Thm.~3]{gretton2005hsic}.
\end{lemma}

\begin{proof}
The estimator form: with $K_Z=ZZ^{\!\top}$,
$\widehat{\HSIC}=(m-1)^{-2}\tr(K_ZHK_yH)=(m-1)^{-2}\tr(Z^{\!\top}HK_yHZ)$ by cyclicity
(and $M=HK_yH$). For the population claim, the linear kernel on $Z$ has feature map
$z\mapsto z$, so the cross-covariance operator factorizes over coordinates and
$\HSIC=\sum_a\|C_{z_a\varphi(y)}\|^2_{\mathrm{HS}}$ with
$\|C_{z_a\varphi(y)}\|_{\mathrm{HS}}=\sup_{\|g\|_{\mathcal H_y}\le1}
\operatorname{Cov}\!\big(z_a,\,g(y)\big)$. Thus $\HSIC=0$ iff
$\operatorname{Cov}(z_a,g(y))=0$ for all $g\in\mathcal H_y$ and all $a$. The Gaussian
RKHS is dense in $L^2(P_y)$ for any distribution $P_y$ on $\R$
\cite{steinwart2008book,sriperumbudur2010}, so the condition extends to all
$g\in L^2(P_y)$; choosing $g(y)=\E[z_a\mid y]$ gives
$\operatorname{Var}\big(\E[z_a\mid y]\big)=0$, i.e.\ $\E[z_a\mid y]$ constant a.s.,
for every $a$. The converse is immediate from the tower rule.
\end{proof}

\begin{remark}
Lemma~\ref{lem:hsic} is the precise sense in which the $\lambda=0$ end is
``supervised'': the probe maximizes dependence of the embedding on the year in
\emph{any functional form} (the Gaussian kernel sees all of $L^2$), but through the
first moment of $Z$. This is exactly the right target for a downstream regression
readout, which consumes $\E[Z\mid y]$-type signal.
\end{remark}

\section{Exactness of the returned solution}

\begin{theorem}[Correctness]\label{thm:main}
For every $\lambda\in[0,1]$, the value of program \eqref{eq:program} equals
$\sum_{j=1}^{d}\gamma_j(\lambda)$, and it is attained by
$V^\ast=[v_1,\dots,v_d]$, the $B$-orthonormal eigenvectors of the pencil
$(A_\lambda,B)$ associated with the $d$ largest eigenvalues. This is precisely the
matrix returned by the algorithm; hence the algorithm outputs an exact global maximizer
of its objective --- no relaxation, no local search.
\end{theorem}

\begin{proof}
By Lemma~\ref{lem:pd}, $B\succ0$; substitute $W=B^{1/2}V$, a bijection between
$\{V:V^{\!\top}BV=I_d\}$ and the Stiefel manifold $\{W:W^{\!\top}W=I_d\}$, under which
$J_\lambda(V)=\tr(W^{\!\top}\tilde A_\lambda W)$ with
$\tilde A_\lambda=B^{-1/2}A_\lambda B^{-1/2}$ symmetric. Let
$\tilde A_\lambda=\sum_{j=1}^{r}\gamma_j u_ju_j^{\!\top}$ be its spectral decomposition,
$\gamma_1\ge\dots\ge\gamma_r$ (these are the pencil eigenvalues: $\tilde A_\lambda w=\gamma w
\iff A_\lambda v=\gamma Bv$ with $v=B^{-1/2}w$). For any feasible $W$,
\[
\tr\!\big(W^{\!\top}\tilde A_\lambda W\big)
=\sum_{j=1}^{r}\gamma_j\,\underbrace{\|W^{\!\top}u_j\|_2^2}_{=:c_j},
\qquad
c_j = u_j^{\!\top}WW^{\!\top}u_j .
\]
Since $WW^{\!\top}$ is an orthogonal projection of rank $d$, we have $0\le c_j\le
\|u_j\|^2=1$ and $\sum_j c_j=\tr(WW^{\!\top})=d$. Hence
$\tr(W^{\!\top}\tilde A_\lambda W)\le\max\{\sum_j\gamma_jc_j : c\in[0,1]^r,\ \sum_jc_j=d\}
=\sum_{j\le d}\gamma_j$, the maximum of a linear functional over that polytope being
attained at the vertex $c=(1,\dots,1,0,\dots,0)$ ordered by $\gamma$. The bound is
achieved by $W^\ast=[u_1,\dots,u_d]$, i.e.\ $V^\ast=B^{-1/2}W^\ast$, which is exactly
the top-$d$ $B$-orthonormal generalized eigenbasis (Ky Fan's maximum principle
\cite{fan1949}, in the generalized-pencil form). Ties among eigenvalues make the
maximizer non-unique but do not affect the optimal value.
\end{proof}

\begin{proposition}[Endpoints]\label{prop:limits}
As $\varepsilon\to0^+$ (and on the column space of $\Xt$):
(i) at $\lambda=1$, program \eqref{eq:program} is equivalent to
$\min_{V^{\!\top}BV=I}\tr(V^{\!\top}\Xt^{\!\top}L\Xt V)$ --- Locality Preserving
Projections \cite{he2003lpp}, the linear out-of-sample form of Laplacian eigenmaps
\cite{belkin2003}; the year never enters.
(ii) at $\lambda=0$ it is supervised PCA with target kernel $K_y$
\cite{barshan2011spca}, computed in the $D$-metric; the graph never enters beyond
the metric. Rescaling $M\mapsto\Mh$, $L\mapsto\Lh$ multiplies the objective at each
endpoint by a positive constant and therefore changes neither maximizer.
\end{proposition}

\begin{proof}
(i) With $\lambda=1$, $J_1(V)=-\tr(V^{\!\top}\Xt^{\!\top}\Lh\Xt V)$, and maximizing $-f$
is minimizing $f$; positive scaling ($\|L\|_2^{-1}$) preserves argmaxima. The
minimization of $\tr(V^{\!\top}\Xt^{\!\top}L\Xt V)$ under $V^{\!\top}\Xt^{\!\top}D\Xt
V=I$ is the LPP program verbatim. (ii) With $\lambda=0$,
$J_0(V)=\tr(V^{\!\top}\Xt^{\!\top}\Mh\Xt V)$, whose maximization under an
orthogonality constraint is the supervised-PCA program of \cite{barshan2011spca} with
$K_y$ as target kernel; our constraint uses the $D$-metric rather than $I$, a
change of inner product that Theorem~\ref{thm:main} handles exactly.
\end{proof}

\begin{proposition}[The constant direction carries nothing and may be discarded]
\label{prop:trivial}
$M\mathbf 1 = 0$ and $L\mathbf 1 = 0$. Consequently, any direction $v$ with embedding
$\Xt v\in\operatorname{span}\{\mathbf 1\}$ has $v^{\!\top}A_\lambda v=0$ for every
$\lambda$: it contributes zero dependence and zero smoothness penalty, i.e.\ zero
information about $y$ and zero geometric content. Removing such directions from the
returned basis (the algorithm's deflation step) changes the objective value by at most
the amount that direction contributed, namely $0$, and is the standard trivial-solution
removal of Laplacian eigenmaps \cite{belkin2003}.
\end{proposition}

\begin{proof}
$H\mathbf 1=0$ gives $M\mathbf 1=HK_yH\mathbf 1=0$; $L\mathbf 1=(D-W)\mathbf 1=0$ since
$d_i=\sum_jw_{ij}$. If $\Xt v=c\mathbf 1$ then
$v^{\!\top}\Xt^{\!\top}\Mh\Xt v=c^2\mathbf 1^{\!\top}\Mh\mathbf 1=0$ and likewise for
$\Lh$.
\end{proof}

\section{Leakage-freeness of the evaluation}

\begin{proposition}[No test-set leakage]\label{prop:leakage}
Fix a draw and a train/test split by ruler groups, $G_{\mathrm{tr}}\cap
G_{\mathrm{te}}=\emptyset$. Every fitted object of the probe --- the PCA map $P$, the
graph $W$ and bandwidth $\sigma_x$, the year kernel $K_y$ and bandwidth $\sigma_y$, the
matrices $(A_\lambda,B)$, the eigenbasis $V^\ast$, and the ridge readout --- is a
measurable function of the training fold alone. Test points enter only through the
pointwise map $x_\ast\mapsto V^{\ast\top}P(x_\ast)$ and the fitted readout. Hence the
reported test statistics are functions of (train fold, test fold) in which no test
observation influences any fitted quantity, and, the split being grouped by ruler, they
estimate generalization to unseen rulers.
\end{proposition}

\begin{proof}
By construction: each listed object is computed from $\{(x_i,y_i):i\in\text{train}\}$
only (the graph is never rebuilt with test nodes; cf.\ the Nystr\"om alternative in
\cite{bengio2004oos}, which shares this property). The composition of a train-measurable
map with a test point is evaluation, not fitting. The group condition makes the test
rulers' identities, and therefore their constant year values, disjoint from training.
\end{proof}

\section{Regularity of the solution path in $\lambda$}

\begin{proposition}[The dial is well behaved]\label{prop:path}
Let $F_d(\lambda)=\sum_{j\le d}\gamma_j(\lambda)$ be the optimal value of
\eqref{eq:program}. Then:
(i) each $\gamma_j(\cdot)$ is Lipschitz on $[0,1]$ with constant
$\big\|\tilde A_M\big\|_2+\big\|\tilde A_L\big\|_2$, where $\tilde A_M =
B^{-1/2}\Xt^{\!\top}\Mh\Xt B^{-1/2}$ and $\tilde A_L$ analogously;
(ii) $F_d$ is convex and continuous on $[0,1]$;
(iii) the spectral projection onto the top-$d$ eigenspace is analytic in $\lambda$
wherever $\gamma_d(\lambda)>\gamma_{d+1}(\lambda)$, so the embedding (up to rotation)
and every continuous readout of it vary continuously in $\lambda$ outside the finitely
many crossing points of these two eigenvalue branches.
\end{proposition}

\begin{proof}
$\tilde A_\lambda=(1-\lambda)\tilde A_M-\lambda\tilde A_L$ is affine in $\lambda$;
(i) is Weyl's perturbation inequality
$|\gamma_j(\lambda)-\gamma_j(\lambda')|\le\|\tilde A_\lambda-\tilde
A_{\lambda'}\|_2=|\lambda-\lambda'|\,\|\tilde A_M+\tilde A_L\|_2
\le|\lambda-\lambda'|(\|\tilde A_M\|_2+\|\tilde A_L\|_2)$.
(ii) By Theorem~\ref{thm:main},
$F_d(\lambda)=\max_{V}\;\tr(V^{\!\top}A_\lambda V)$ over the fixed feasible set; for
each fixed $V$ the map $\lambda\mapsto\tr(V^{\!\top}A_\lambda V)$ is affine, and a
pointwise supremum of affine functions is convex; convex functions on an interval are
continuous in its interior, and (i) supplies continuity at the endpoints.
(iii) is Kato's analytic perturbation theory for symmetric pencils with isolated
eigenvalue groups \cite{kato1966}; eigenvalue branches of a real-analytic (here affine)
symmetric family are analytic, and crossings of the finitely many branch pairs are
isolated unless the branches coincide identically, in which case the projection extends
analytically anyway.
\end{proof}

\begin{remark}[Scope of the guarantee]
Theorem~\ref{thm:main} and Propositions~\ref{prop:limits}--\ref{prop:path} are exact,
finite-sample statements about the \emph{empirical} objective: the algorithm returns the
global optimum of \eqref{eq:program}, interpolates the two named methods at its
endpoints, evaluates without leakage, and traces a continuous path. What they do
\emph{not} assert is statistical fidelity of the empirical objective to its population
counterpart; those are the standard, separately-proved facts that the graph Laplacian
converges to the Laplace--Beltrami operator of the data manifold
\cite{belkin2008towards} and that empirical HSIC concentrates around population HSIC at
rate $O(m^{-1/2})$ \cite{gretton2005hsic}, both under their respective sampling
assumptions. The probe's inferential use in the thesis --- comparing the $\lambda$-curve
of a trained model against a random-initialized control under an identical pipeline ---
requires only the finite-sample guarantees proved here, since both arms share the same
empirical objective and protocol.
\end{remark}

\begin{thebibliography}{9}\small
\bibitem{gretton2005hsic} A.~Gretton, O.~Bousquet, A.~Smola, B.~Sch\"olkopf.
Measuring statistical dependence with Hilbert--Schmidt norms. \emph{ALT}, 2005.
\bibitem{belkin2003} M.~Belkin, P.~Niyogi. Laplacian eigenmaps for dimensionality
reduction and data representation. \emph{Neural Computation} 15(6), 2003.
\bibitem{belkin2008towards} M.~Belkin, P.~Niyogi. Towards a theoretical foundation for
Laplacian-based manifold methods. \emph{JCSS} 74(8), 2008.
\bibitem{he2003lpp} X.~He, P.~Niyogi. Locality preserving projections. \emph{NeurIPS},
2003.
\bibitem{barshan2011spca} E.~Barshan, A.~Ghodsi, Z.~Azimifar, M.~Zolghadri Jahromi.
Supervised principal component analysis. \emph{Pattern Recognition} 44(7), 2011.
\bibitem{fan1949} K.~Fan. On a theorem of Weyl concerning eigenvalues of linear
transformations~I. \emph{PNAS} 35, 1949.
\bibitem{bengio2004oos} Y.~Bengio et al. Out-of-sample extensions for LLE, Isomap, MDS,
eigenmaps, and spectral clustering. \emph{NeurIPS}, 2004.
\bibitem{sriperumbudur2010} B.~Sriperumbudur, A.~Gretton, K.~Fukumizu, B.~Sch\"olkopf,
G.~Lanckriet. Hilbert space embeddings and metrics on probability measures.
\emph{JMLR} 11, 2010.
\bibitem{steinwart2008book} I.~Steinwart, A.~Christmann. \emph{Support Vector
Machines}. Springer, 2008. (Universality/denseness of the Gaussian RKHS.)
\bibitem{kato1966} T.~Kato. \emph{Perturbation Theory for Linear Operators}. Springer,
1966.
\end{thebibliography}

\end{document}
